\documentclass[a4paper]{article}
\usepackage{amsfonts}
\usepackage{amsmath}
\usepackage{amssymb}
\usepackage{graphicx}     

\begin{document}
  
\noindent \large Auteur: Noran Ben Said \\
\noindent \large Studierichting: Informatica\\
\noindent \large Oefeningenreeks 5A nr. 9.6 \\
\medskip

\normalsize

$ I = \displaystyle\int_1^5 \dfrac{x^7}{(6-x)^7 + x^7}\,dx$\\

Stel $ x = 6 - t \Rightarrow dx = -dt$\\

Grenzen:\\

$x=1 \Rightarrow t=5$\\

$x=5 \Rightarrow t=1$\\

$ I = \displaystyle\int_5^1 \dfrac{(6-t)^7}{t^7 + (6-t)^7}(-dt)$\\

$ = \displaystyle\int_1^5 \dfrac{(6-x)^7}{x^7 + (6-x)^7}\,dx$\\

$ I + I = \displaystyle\int_1^5 \dfrac{x^7 + (6-x)^7}{x^7 + (6-x)^7}\,dx$\\

$ 2I = \displaystyle\int_1^5 1\,dx$\\

$ = \left[x\right]_1^5  = 5 - 1 = 4$\\

$ I = 2$\\


\begin{thebibliography}{999}
%\bibitem{a} Referentie
\end{thebibliography}
\end{document} 


